\documentclass{article}
\usepackage{amsmath,amssymb,url}
\title{ARC13: Exact order-independent products and random-access audit}
\date{2026-10-09}
\begin{document}\maketitle
\paragraph{Status and corpus.}
COMPUTATIONAL candidate; all-input favorable compilation and P versus NP OPEN.
Adopt recovered replay modules and ARC11--12 from this conversation.
ARC12 handles low-run sequences; ARC13 avoids sorting entirely for
multiset equality. Root main.tex ownership is preserved.
This is not an integrated recursive SAT solver. Conditional K elimination
is unchanged. No retired cosine endpoint is used.

\paragraph{Exact construction.}
For any two residual clause multisets A,B with the same count m,
fix an injective code into the field of prime order p, with p>m.
Normalize literals using fixed labels, delete tautologies and satisfied
clauses, preserve duplicate clauses and the empty clause. Define
\[
 P_A(z)=\prod_{c\in A}(z-\operatorname{code}(c)).
\]
Compare products at m+1 distinct field elements. Equality at every point
implies P_A=P_B since their difference has degree at most m.
Unique factorization into monic linear factors gives A=B as multisets.
Counts differing or any product mismatch gives DISTINCT as multisets,
not necessarily different SAT answers. Equal multisets suffice for SAT reuse.
A matching capped prefix gives UNKNOWN. This is deterministic, not a
probabilistic hash acceptance rule. Equal-degree monicity permits tightening
the point count, but the implementation conservatively uses m+1.

\paragraph{Resource ledger.}
Only a current normalized clause, products, counters, field description,
source cursor and fixed assignments are needed. For bounded width, polynomial
range codes and a polynomial-size field, logical workspace is logarithmic
plus assignment storage. The implemented repeated-evaluation scheme performs
O(mN) source work plus charged modular arithmetic; discovery count scans and
failed evaluations are included. Repeated invocations, child executions,
ancestor assignments and field acquisition remain charged in composition.
Selecting a suitable prime uniformly is not free. The finite executable
uses p=65537, verifies primality by trial division, rejects out-of-field
clause codes, and rejects degree at least p. This is a finite-domain
experiment, not a uniform large-input field compiler.

\paragraph{Validation and falsification.}
256 exhaustive pairs of clause subsets pass exact multiset comparison.
Eight shuffled tests at sizes 8,32,128,512 use seed 13 and caps 4 or m+1.
At 512 clauses complete equality requires 513 evaluation pairs and 2105344
source literal reads; cap 4 returns UNKNOWN after 20480 reads.
These scaling lists contain repeated positive unit clauses and are easy
SAT: costly certification is not evidence of SAT hardness.
A single evaluation is insufficient: multisets [1,6] and [2,3] both have
product 6 at zero but different products at one.
The executable finishes below 25 seconds, with an internal 24-second alarm.
An independent agent confirmed the theorem and independently enumerated
small-field multisets. Tests are not all-input complexity proofs.

\paragraph{Random-access lower bounds: what transfers.}
Beame, A General Sequential Time-Space Tradeoff for Finding Unique Elements,
SIAM Journal on Computing, doi:10.1137/0220017, proves an
Omega(m squared) time-space product for the stated unique-elements output
problem and sorting on R-way branching programs. The publisher abstract
explicitly includes offline log-cost RAMs among covered models.
Thus random access alone does not make general sorting cheap.
However this is not a direct lower bound for Boolean multiset equality,
clause-set equality, or SAT. Importing it requires an explicit
resource-preserving reduction. Our comparator avoids sorted output.
No such SAT reduction is proved here.

\paragraph{Verified SAT endpoint.}
Fortnow, Lipton, van Melkebeek and Viglas, Time-Space Lower Bounds for
Satisfiability, primary author manuscript dated September 7, 2004,
\url{https://lance.fortnow.com/papers/files/tst.pdf},
Section 2.1 uses deterministic random-access Turing machines with charged
work and index tapes. Theorem 3.1 gives a locally computable SAT reduction
with polylogarithmic time overhead and logarithmic auxiliary space.
Corollary 5.8 gives, for c>1,
\[
 \operatorname{NTIME}[n]\not\subseteq
 \operatorname{DTISP}[n^c,o(n^{d(c)})],\qquad
 d(c)=\tfrac12((c+2)/c^2-c).
\]
Choosing c=1201/1000 gives
d(c)=1468676399/2884802000>1/2. An actual uniform all-length SAT algorithm
with simultaneous N to the power (6/5+o(1)) time and
N to the power (1/2+o(1)) work space in the matched model would therefore
contradict this endpoint after absorbing polylogarithmic overhead.
Neither ARC13 nor ARC12 supplies that algorithm.
Random access is already included, not a loophole or a new separation.

\paragraph{Constructive literature clue and next target.}
Theorem 3.1 explicitly avoids storing a reduction or recomputing it wholesale:
each requested output bit is computed in polylogarithmic time and logarithmic
workspace. Recover this locally addressable pattern for wrapper descendants.
Its known reduction is quasi-linear in size, not a SAT shrinker; the missing
bridge is simultaneous shrinkage, cheap bit access and charged coverage.
Prefer proving such an accessor over assuming sorting must be performed.
\paragraph{Sentinel.} ARC13-END-OPEN.
\end{document}

